\color{unidarkgreen}{\subsection{Орган выявления броска тока намагничивания компенсационной обмотки}}
\color{black}

\begin{enumerate}

\item
Для обеспечения несрабатывания «МТЗ КО», «ТЗНП КО», «ТЗОП КО» при броске тока намагничивания используется блокировка по содержанию 2 гармоники в токе. Блокировка основана на контроле отношения действующего значения второй гармонике к первой в токах фаз. Имеется возможность регулировать порог срабатывания блокировки, с помощью уставки \textbf{«I2г/I1г БНТ»}. Блокировка по 2 гармонике выводится из действия при величине тока ниже 0,04Iном или превышении током первой гармоники значения \textbf{«Iср БНТ»}.  

\item
Алгоритм работы ОВ БНТ КО приведен на рисунке \ref{pict_bntko:BNT_KO}. Уставки ОВ БНТ КО приведены в таблице \ref{app_set:bnt}.

\medskip
\begin{figure}[H]
\centering
\includegraphics[width=0.9\textwidth,height=0.9\textheight,keepaspectratio]{\fbpath/11.02 УШР/ОВ БНТ КО/img/BNT_KO.pdf}
\captionof{figure}{Функциональная схема логики ОВ БНТ КО}\label{pict_bntko:BNT_KO}
\end{figure}



\end{enumerate}